\documentclass{article}
\usepackage{verbatim}
\usepackage{hyperref}
\usepackage{fullpage}
\usepackage{amsmath}
\usepackage{amssymb}
\usepackage{graphicx}
\usepackage[shortlabels]{enumitem}

\title{Problem Set 1}
\begin{document}
\small
\date{}
\author{Prof. Conlon}
\maketitle
Your answers should be produced in \LaTeX, and should include all relevant graph and code.  Code should be in the appropriate verbatim environment and properly documented. You are allowed to work in groups but you must turn in your own writeup. Submit your assignment via Brightspace.

You may want to consult my bonus notes on nonlinear optimization and/or GMM.\\

Consider the following utility model for parents choosing schools
\begin{align*}
u_{ij} = \beta_1 \cdot \text{test scores}_j  +\beta_2 \cdot \text{sports}_j+ \xi_j - \alpha \cdot d_{ij} +  \varepsilon_{ij}
\end{align*}
We define the terms below:
\begin{itemize}
	\item Test Scores $\in [0,1]$: the average percentile (as a fraction) of the students at the school on a standardized test
	\item Sports $\in [0.3, 4]$ (roughly): how good the school is at sports
	\item Distance $d_{ij} \in [0,\infty)$ how far away the school is in miles from household $i$
	\item $\xi_j$: unobservable school ``quality''
\end{itemize}
There is no outside option: every household chooses one of the five schools. That means only \emph{differences} in the $\xi_j$'s are identified, so normalize $\xi_1 = 0$. The full parameter list is $\theta = [\beta_1,\beta_2, \xi_2,\ldots,\xi_5, \alpha]$ --- though as you will discover below, you cannot estimate all of these at once.

\section*{\normalsize Questions}
For parameter estimates, also report: the (negative) log-likelihood, the diversion ratio from school 1 to school 2, the average own elasticity with respect to distance for the chosen school.
\begin{enumerate}
	\item Look at the data and plot the distribution of distance to all schools, and the distribution of distance to the chosen school.
	\item Write down the choice probability $\sigma_{ij}$ and log-likelihood for a plain logit model. (This is micro data: each household has its own $\sigma_{ij}$; save $s_j$ for observed shares.)
	\item Write down the score and the gradient of your log-likelihood.
	\item Estimate the \emph{characteristics} model by maximum likelihood: impose $\xi_j = 0$ for all $j$, so that $\theta = [\beta_1,\beta_2,\alpha]$.
	\item Report standard errors for your estimates using the inverse Hessian (observed information). Compare with standard errors from the outer product of scores.
	\item Estimate the \emph{fixed effects} model with only the $\xi_j$'s and $\alpha$ (i.e., impose $\beta_1=\beta_2=0$, so that $\theta = [\xi_2,\ldots,\xi_5,\alpha]$). Add that to your table. Show that the characteristics model from Question 4 is \emph{nested} in this one: what restriction on the $\xi_j$'s does it impose? Compare the two with a likelihood-ratio test --- how many degrees of freedom?
	\item \label{q:ident} Why can't you estimate $\beta_1$, $\beta_2$, and all of the $\xi_j$'s jointly? (What varies \emph{within} a school across households, and what doesn't?) This is not a small-sample problem: it is exactly the situation in BLP-style demand models, where every product carries an unobservable $\xi_j$ alongside its characteristics $x_j$, and it is why product fixed effects ``absorb'' product characteristics. Now run the two-step fix: regress your estimated $\widehat{\xi}_j$'s from Question 6 on test scores and sports. How do the coefficients compare to Question 4's $\widehat{\beta}$'s, and why? With $J=5$ schools, would you trust the standard errors from that regression? (We will see this move again when we get to BLP.)
	\item Now allow for parents to have different preferences for $ \text{test scores}_j$ so that $\beta_{1i}\sim N(\beta_1,\Sigma_b^2)$, where the standard deviation $\Sigma_b$ is just a scalar here. (Build on the characteristics model from Question 4, keeping $\xi_j = 0$.) Write down the (simulated) choice probability and gradient expressions.

	\item \label{q:rcml} Estimate this expanded model via maximum likelihood:
	\begin{enumerate}
		\item Using 100 Monte Carlo Draws from an appropriately transformed standard normal.
		\item Using a Gauss Hermite quadrature rule.
	\end{enumerate}
	\item Read Chapter 10 in Train and write down the MSM estimator for the expanded model. What are your ``instruments''?
	\item Calculate the Jacobian of the MSM estimator.
	\item Estimate the Parameters of the MSM estimator.
	\item Report standard errors using the GMM sandwich formula. What adjustment does simulation error require? (Hint: with $S$ simulation draws the variance picks up a $(1+1/S)$ factor.)
	\item Bonus: Using your initial MSM estimates as a starting point, explain how to construct an ``efficient'' MSM estimator, and produce ``efficient'' estimates.
	\item Suppose school 1 closes. Using your mixed logit estimates from Question \ref{q:rcml}, compute each household's change in consumer surplus via the Small--Rosen log-sum-exp difference:
	\begin{align*}
	\Delta CS_i = \frac{1}{\alpha} \, \mathbb{E}_{\beta_{1i}}\left[ \log \sum_{j \neq 1} \exp(V_{ij}) - \log \sum_{j=1}^{5} \exp(V_{ij}) \right]
	\end{align*}
	Note that distance plays the role of price here, so the ``money'' metric is in miles of extra travel. Report the average $\Delta CS$ and a histogram across households. Compare with the plain-logit shortcut $\Delta CS_i = \frac{1}{\alpha} \log(1-\sigma_{i1})$ evaluated at your plain logit estimates. Why do the two differ? (See the Welfare frames from the IIA Logit lecture.)
\end{enumerate}

\end{document}

